% ============================================================
%  第 5 章  安全设计与实现
% ============================================================
\section{安全设计与实现}
\label{sec:security}

本章论述系统自身的安全设计。需要区分两个概念：\cref{sec:algorithms} 讨论的是
系统\emph{如何检测攻击}，本章讨论的是系统\emph{自身如何不成为攻击面}。
安全产品若自身存在缺陷，其危害往往甚于被防护的威胁。

\subsection{纵深防御体系}

系统对同一类威胁在多个层次分别施加约束，任一层失效不导致整体防护失守。
\cref{tab:defense-depth} 以 SQL 注入为例展示这一分层结构。

\begin{table}[htbp]
\centering
\caption{针对 SQL 注入的纵深防御层次}
\label{tab:defense-depth}
\small
\begin{tabular}{@{}c l l l@{}}
\toprule
\textbf{层次} & \textbf{位置} & \textbf{防护手段} & \textbf{失效后果} \\
\midrule
第一层 & 网关 & 参数攻击意图启发式判定 & 明显攻击未被提前滤除 \\
第二层 & 探针 & ��实语句的结构指纹比对 & 失去确定性拦截能力 \\
第三层 & 应用 & 参数化查询与输入类型校验 & 注入得以实际执行 \\
第四层 & 数据库 & 最小权限账户与视图隔离 & 数据泄露范围扩大 \\
\bottomrule
\end{tabular}
\end{table}

\begin{remark}
第三层的参数化查询是\emph{根本性}修复：它使 SQL 的结构与数据彻底分离，
注入在原理上不可能发生。前两层属于\emph{补偿性}控制，其价值在于覆盖那些
因存量系统改造成本、第三方组件不可控等现实约束而无法��用第三层修复的场景。
本文实现的检测能力应被理解为纵深防御中的一环，而非参数化查询的替代品。
\end{remark}

\subsection{认证与授权}

\subsubsection{令牌签名算法的强制约束}

JSON Web Token 规范允许 \code{alg} 字段取值为 \code{none}，表示无签名。
若校验方未\emph{强制指定}期望算法，攻击者可将令牌头部改为
\code{\{"alg":"none"\}}、删除签名段，并任意篡改载荷中的身份与角色声明。
这构成威胁 T-03，CVSS 评分 9.8。

正确的做法是：\textbf{信任服务端配置的算法，绝不信任令牌自称的算法}。
本文在令牌解析时显式绑定密钥与签名算法，使解析器拒绝任何算法不匹配的令牌，
由此同时防御无签名令牌与算法降级两类攻击。

\subsubsection{认证与授权的分离}

一个常见的设计缺陷是将\zh{已登录}等同于\zh{有权限}。本文明确区分二者：
认证回答\zh{你是谁}，授权回答\zh{你能做什么}。管理接口在验证令牌有效性之外，
额外校验角色声明；资源访问接口则强制以\emph{服务端会话中的主体标识}
作为查询条件，而非信任客户端传入的标识。

后一点尤为关键。水平越权（不安全的直接对象引用）的根源正是：
\emph{权限判定的依据来自用户可控的输入}。将属主标识从查询参数改为
会话属性，即从设计上消除了该类漏洞的可能。

\subsection{敏感数据保护}

\subsubsection{口令存储}

系统采用 BCrypt 自适应哈希算法存储口令凭据\cite{provos2003bcrypt}，
工作因子设为 12。相较常见的无盐 MD5 方案，其改进体现在两个维度：

\begin{itemize}
  \item \textbf{内置随机盐}：相同口令每次哈希产生不同结果，
        使预计算的彩虹表失效。
  \item \textbf{计算代价可调}：工作因子可随硬件性能提升而调高，
        使暴力破解的时间成本始终维持在不可接受的水平。
\end{itemize}

\subsubsection{出站请求的地址约束}

服务端请求伪造（威胁 T-10）的防护有一处易被忽视的细节：
\textbf{必须先解析域名得到真实地址再做判定}，而非仅检查 URL 文本。
这一步同时防御两类绕过：进制编码绕过（\code{http://2130706433/} 与
\code{http://0177.0.0.1/} 解析后均指向回环地址）与域名重绑定
（攻击者控制的域名可解析至内网地址）。判定范围覆盖回环地址、私有网段、
链路本地地址与云平台元数据端点。

\subsection{审计日志的不可否认性}
\label{subsec:hashchain}

\subsubsection{哈希链设计}

审计日志若可被篡改，则失去取证价值。攻击者在得手后抹除日志痕迹，
是入侵活动的标准环节。为此本文采用\emph{哈希链}结构，使日志具备
防篡改的密码学保证。

设第 $i$ 条事件记录的内容摘要为 $m_i$，则其哈希值定义为
\begin{equation}
h_i = H(m_i \,\Vert\, h_{i-1}), \qquad h_0 = \varepsilon
\label{eq:hashchain}
\end{equation}
其中 $H$ 为 SHA-256 哈希函数。每条记录保存 $h_{i-1}$ 与 $h_i$ 两个字段，
形成链式依赖，如 \cref{fig:hashchain} 所示。

\begin{figure}[htbp]
\centering
\begin{tikzpicture}[node distance=6mm]
\node[block, minimum width=25mm, minimum height=14mm] (e1)
  {事件 $E_1$\\[-2pt]{\scriptsize $h_0 = \varepsilon$}\\[-2pt]
   {\scriptsize $h_1 = H(m_1 \Vert h_0)$}};
\node[block, right=7mm of e1, minimum width=25mm, minimum height=14mm] (e2)
  {事件 $E_2$\\[-2pt]{\scriptsize $\mathrm{prev} = h_1$}\\[-2pt]
   {\scriptsize $h_2 = H(m_2 \Vert h_1)$}};
\node[blockred, right=7mm of e2, minimum width=27mm, minimum height=14mm] (e3)
  {事件 $E_3$（被篡改）\\[-2pt]{\scriptsize $\mathrm{prev} = h_2$}\\[-2pt]
   {\scriptsize $h_3' \neq H(m_3' \Vert h_2)$}};
\node[block, right=7mm of e3, minimum width=25mm, minimum height=14mm] (e4)
  {事件 $E_4$\\[-2pt]{\scriptsize $\mathrm{prev} = h_3$}\\[-2pt]
   {\scriptsize 校验失败}};

\draw[flow] (e1) -- (e2);
\draw[flow] (e2) -- (e3);
\draw[flowred] (e3) -- (e4);

\node[below=3mm of e2, align=center, lbl, text width=88mm, xshift=22mm]
  {任一记录被篡改，其自身哈希与后续所有记录的前置哈希校验均失败，篡改行为无法隐藏};
\end{tikzpicture}
\caption{哈希链式审计日志的防篡改机制}
\label{fig:hashchain}
\end{figure}

\begin{proposition}[篡改可检测性]
在 $H$ 抗第二原像的前提下，对已存储的任一记录 $E_i$ 的任何修改，
均将导致完整性校验在位置 $i$ 或其后失败。
\end{proposition}

\begin{proof}
设攻击者将 $m_i$ 篡改为 $m_i' \neq m_i$。若其不修改 $h_i$，则重算得
$H(m_i' \Vert h_{i-1}) \neq h_i$，校验在位置 $i$ 失败。若其同时将 $h_i$
更新为 $h_i' = H(m_i' \Vert h_{i-1})$，则记录 $E_{i+1}$ 中存储的前置哈希
仍为 $h_i \neq h_i'$，校验在位置 $i+1$ 失败。归纳可知，攻击者须重算
从 $i$ 至链尾的全部哈希方能保持一致，而这在链尾哈希被独立留存或
定期公证的前提下同样可被检出。
\end{proof}

\subsubsection{实现中的一处关键细节}

哈希链的实现有一个易犯的错误：\emph{清空日志时未同步重置内存中缓存的
链尾哈希}。若仅删除数据库记录而保留旧的链尾哈希，后续写入的第一条记录
将携带指向已删除记录的前置哈希，导致链条从起点即断裂，完整性校验永远失败。

本文在开发过程中确实遇到了这一问题，并将存储清空与链状态重置封装为
单一的原子操作予以修复。此类\zh{持久化状态与内存状态不同步}的缺陷
在各类重置操作中具有普遍性，值得在设计阶段即予以关注。

\subsection{可用性保障}

\subsubsection{失效安全策略}

检测组件的任何内部异常均降级为放行并记录，而非阻断业务。这一取舍看似
削弱了安全性，实则是经过权衡的工程决策：安全组件的可用性故障若转化为
全站业务中断，其造成的损失通常超过单次攻击得手。

实现上需注意一处细节：\emph{拦截决策本身不能被失效安全逻辑吞掉}。
探针的检测方法在捕获异常时需区分两类情形——检测过程的意外异常应被吞掉
并放行，而主动抛出的安全拦截异常必须向上传播以阻断执行。本文在实现中
曾因字节码注入框架的异常抑制选项配置不当而导致\zh{能检出却无法拦截}，
后经修正。

\subsubsection{速率控制}

系统采用令牌桶与滑动窗口双重限流。二者的组合并非冗余，而是针对不同的
攻击模式：令牌桶控制\emph{瞬时突发}，其桶容量决定了可容忍的突发规模；
滑动窗口控制\emph{持续总量}，消除了固定窗口计数在窗口边界处可通过双倍
请求的缺陷。两者取合取关系，任一超限即拒绝。

\subsubsection{来源信誉的时间衰减}

单次攻击可能是误报，但同一来源的持续攻击是明确的恶意信号。系统为每个
来源地址维护累积风险评分，达到阈值后临时封禁。评分按下式随时间衰减：
\begin{equation}
s(t) = s_0 \cdot \gamma^{\,\Delta t}, \qquad \gamma = 0.9
\end{equation}
其中 $\Delta t$ 为距上次事件的小时数。

衰减机制对降低误报至关重要。共享出口地址（如校园网地址转换出口）
可能因个别用户的行为被误判，衰减保证正常用户不会被长期阻断。
同理，封禁设有明确时限而非永久，符合失效安全原则。

\begin{remark}[阈值设定的实践考量]
封禁阈值的设定需权衡两种失效：阈值过低时，安全测试或演示场景下来自同一
来源的连续请求会很快触发封禁，后续所有请求被统一拦下，反而\emph{掩盖了
各类攻击的具体检测结果}；阈值过高则削弱对真实持续攻击的抑制能力。
本文在测试中观察到前一种情形，据此将阈值调整至相当于连续十次严重攻击的水平。
\end{remark}

\subsection{安全编码实践}

\subsubsection{高危函数禁用}

依据课程实践指导书的要求，本文明确禁用一组高危函数并给出替代方案，
如 \cref{tab:banned-func} 所示。

\begin{table}[htbp]
\centering
\caption{高危函数禁用清单与替代方案}
\label{tab:banned-func}
\small
\begin{tabular}{@{}l l p{4.4cm}@{}}
\toprule
\textbf{禁用项} & \textbf{替代方案} & \textbf{原理说明} \\
\midrule
字符串拼接构造 SQL & 参数化查询 & 结构与数据分离，注入原理上不可能 \\
单字符串形式的命令执行 & 参数数组化调用 & 不经 shell 解析，元字符失去语义 \\
原生对象反序列化 & JSON 加类型白名单 & 消除任意类实例化的可能 \\
无盐 MD5 存储口令 & BCrypt 自适应哈希 & 随机盐加可调计算代价 \\
正则过滤跨站脚本 & 文档对象模型白名单净化 & 语义判定优于文本匹配 \\
普通随机数生成令牌 & 密码学安全随机源 & 保证不可预测性 \\
\bottomrule
\end{tabular}
\end{table}

\subsubsection{安全注释规范}

为便于安全审查与代码检索，本文对安全相关代码建立统一的标记约定：
\code{[SEC-xxx-nn]} 标识安全防护逻辑，\code{[VULN-nn]} 标识受控漏洞，
\code{[FIX-nn]} 标识对应的修复实现。每处安全逻辑的注释须说明其
\emph{安全原理}与\emph{对应威胁编号}，而非仅描述代码行为。
\cref{lst:sec-comment} 给出示例。

\begin{lstlisting}[style=javastyle, caption={安全注释规范示例}, label={lst:sec-comment}]
/**
 * [SEC-SQLI-03] AST 字面量归一化
 *
 * 安全原理：
 *   遍历语法树，将所有字面量替换为占位符，保留表名、列名、函数名
 *   与运算符等决定语义结构的节点。归一化后的语句仅反映 SQL 的
 *   骨架形状，与具体参数取值无关。
 *
 *   这是结构指纹能够抵御编码混淆的根本原因：攻击者可任意改变参数
 *   的文本形态，但只要想改变查询逻辑（这正是注入的目的），就必须
 *   引入新的语法节点，骨架必然改变。
 *
 * 威胁对应：T-01 (STRIDE: Tampering, DREAD 9.0)
 */
public static AstNode normalize(AstNode root) { ... }
\end{lstlisting}

\subsubsection{受控漏洞的隔离}

靶场中的漏洞代码为教学目的主动构造，须严格隔离以防误用。本文采取三项措施：
全部漏洞实现集中于独立的 \code{vuln} 包；每个漏洞方法附加显式警示注释，
说明其危害与正确做法；每个漏洞均在 \code{safe} 包提供对应的安全实现。
后者的设计使得\zh{漏洞修复证明}直接体现为\emph{同一功能的两份代码对照}，
既便于教学演示，也为 \cref{sec:evaluation} 的第三轮测试提供了实验基础。

\subsection{系统自身的攻击面分析}

安全产品自身可能成为攻击目标。本文识别并处置了三处潜在风险。

\paragraph{事件上报接口的伪造。}
若上报接口无认证，攻击者可注入大量伪造事件以制造告警噪声，
或伪造\zh{已拦截}记录掩盖真实攻击。本文对内部上报接口采用共享密钥认证。

\paragraph{检测逻辑的递归触发。}
探针钩取数据库执行点后，其自身在上报事件或查询基线时产生的数据库操作
将再次触发检测，形成无限递归。本文通过线程局部标记识别\zh{正在执行检测}
的状态并跳过，避免该问题。

\paragraph{解析资源耗尽。}
攻击者可构造超长载荷或深层嵌套结构，使解析过程消耗大量计算资源。
本文对载荷长度、解码深度与递归深度均设置上限，超限即截断并计入风险信号。
